
\section{Постановка задачи}

Задан линейный стационарный объект
\[
\dot{\bm x}=A\bm x+\bm b u
\]

Матрицы системы и параметры критерия качества
\[
A=\begin{pmatrix}
0&1\\
-2&9
\end{pmatrix},
\qquad
\bm b=\begin{pmatrix}4\\8\end{pmatrix},
\qquad
Q=\begin{pmatrix}3&0\\0&2\end{pmatrix},
\qquad r=1
\]

Критерий качества на конечном интервале времени
\[
J=\int_0^5
\left(
\bm x^T(\tau)Q\bm x(\tau)+ru^2(\tau)
\right)d\tau
\]

Начальные условия
\[
\bm x(0)=\begin{pmatrix}1\\0\end{pmatrix}
\]

\section{Расчёт оптимального регулятора методом Беллмана}

Функцию Беллмана представим в квадратичном виде
\[
V(t,\bm x)=\bm x^TP(t)\bm x
\]

Уравнение Гамильтона-Якоби-Беллмана имеет вид
\[
-\frac{\partial V}{\partial t}
=
\min_u
\left(
\bm x^TQ\bm x+ru^2+
\frac{\partial V}{\partial\bm x}^T
(A\bm x+\bm b u)
\right)
\]

Вычислим производные
\[
\frac{\partial V}{\partial t}
=\bm x^T\dot P\bm x,
\qquad
\frac{\partial V}{\partial\bm x}=2P\bm x
\]

Подставляя их в уравнение Беллмана, получаем
\[
\begin{aligned}
0={}&
\bm x^T(\dot P+A^TP+PA+Q)\bm x+\min_u
\left(ru^2+2u\bm b^TP\bm x\right)
\end{aligned}
\]

Условие минимума по управлению
\[
2ru+2\bm b^TP\bm x=0
\]

Отсюда
\[
u^*(t)=-r^{-1}\bm b^TP(t)\bm x(t)
\]

Подставляя оптимальное управление, получаем
дифференциальное уравнение Риккати
\[
-\dot P=A^TP+PA+Q-P\bm b r^{-1}\bm b^TP
\]

Поскольку терминальный штраф отсутствует,
конечное условие имеет вид
\[
P(5)=0
\]

Матрицу $P(t)$ представим в симметричном виде
\[
P(t)=
\begin{pmatrix}
p_{11}(t)&p_{12}(t)\\
p_{12}(t)&p_{22}(t)
\end{pmatrix}
\]

Подставляя заданные матрицы в уравнение Риккати,
получаем систему дифференциальных уравнений
\[
\begin{cases}
-\dot p_{11}=-4p_{12}+3-
(4p_{11}+8p_{12})^2\\[2mm]
-\dot p_{12}=p_{11}+9p_{12}-2p_{22}-
(4p_{11}+8p_{12})(4p_{12}+8p_{22})\\[2mm]
-\dot p_{22}=2p_{12}+18p_{22}+2-
(4p_{12}+8p_{22})^2
\end{cases}
\]

С конечными условиями
\[
p_{11}(5)=p_{12}(5)=p_{22}(5)=0
\]

Численно интегрируя систему от $t=5$ к $t=0$,
получаем
\[
P(0)\approx
\begin{pmatrix}
2{,}676107&-1{,}731905\\
-1{,}731905&1{,}496937
\end{pmatrix}
\]

Коэффициенты оптимального регулятора зависят от времени
\[
K(t)=r^{-1}\bm b^TP(t)
\]

С учётом найденной матрицы $P(t)$ получаем
\[
K(t)=
\begin{pmatrix}
4p_{11}(t)+8p_{12}(t)&
4p_{12}(t)+8p_{22}(t)
\end{pmatrix}
\]

Таким образом, оптимальное управление имеет вид
\[
\begin{aligned}
u^*(t)
={}&-\left(4p_{11}(t)+8p_{12}(t)\right)x_1(t)\\
&-\left(4p_{12}(t)+8p_{22}(t)\right)x_2(t)
\end{aligned}
\]

В начальный момент времени
\[
u^*(0)=3{,}150813x_1(0)-5{,}047877x_2(0)
\]

К концу интервала коэффициенты регулятора
стремятся к нулю вследствие конечного условия
$P(5)=0$

\section{Моделирование оптимального управления}

Замкнутая система с оптимальным регулятором
описывается уравнением
\[
\dot{\bm x}=
\left(A-\bm bK(t)\right)\bm x
\]

Критерий качества вычисляется по формуле
\[
J(t)=\int_0^t
\left(
3x_1^2(\tau)+2x_2^2(\tau)+u^2(\tau)
\right)d\tau
\]

Минимальное значение критерия качества
определяется функцией Беллмана
\[
J^*=V(0,\bm x(0))
=\bm x^T(0)P(0)\bm x(0)
\]

С учётом начальных условий получаем
\[
\begin{aligned}
J^*
&=
\begin{pmatrix}1&0\end{pmatrix}
P(0)
\begin{pmatrix}1\\0\end{pmatrix}\\
&=p_{11}(0)\approx2{,}676107
\end{aligned}
\]

\section{Исследование отклонений параметров регулятора}

Рассмотрим изменение коэффициентов оптимального
регулятора на $10\%$ и $20\%$
\[
K_\alpha(t)=\alpha K(t)
\]

Выберем 5 значений параметра
\[
\alpha_1=0{,}8,\quad
\alpha_2=0{,}9,\quad
\alpha_3=1,\quad
\alpha_4=1{,}1,\quad
\alpha_5=1{,}2
\]

Соответствующие законы управления
\[
u_\alpha(t)=-\alpha K(t)\bm x(t)
\]

Таким образом, отклонение коэффициентов регулятора
от оптимальных значений приводит к увеличению
критерия качества, что подтверждает оптимальность
управления, полученного методом Беллмана.
